\documentclass[review]{elsarticle}

\usepackage[T1]{fontenc} % agar karakter < > tercetak benar di mode teks
\usepackage{amssymb}
\usepackage{amsmath}
\usepackage{lineno}
\usepackage{hyperref}
\usepackage{graphicx}  % untuk figure
\usepackage{array}     % kolom p{} rata tengah (>{\centering\arraybackslash})
\usepackage{seqsplit}  % memecah string panjang tanpa spasi (hash JA3/JA4)
\usepackage{listings}  % code block dengan auto line-wrap
\usepackage{xcolor}
\lstset{
  basicstyle=\ttfamily\small,
  breaklines=true,
  breakatwhitespace=false,
  columns=fullflexible,
  frame=single,
  backgroundcolor=\color{gray!5},
  showstringspaces=false
}
\setlength{\tabcolsep}{4pt} % tabel lebih rapat agar muat di lebar halaman
\modulolinenumbers[5]

\journal{Computers \& Security}

\begin{document}

\begin{frontmatter}

\title{Adaptive Evasion Traffic Generator (AETG): A Framework for Adaptive IDS Validation via Red-Team Techniques}

\author{Afiq Andico Pangimpian}
\ead{afiqandico13@gmail.com}

\address{Institut Teknologi dan Bisnis STIKOM Bali, Indonesia}

\begin{abstract}
The increasing adoption of machine learning (ML) for network intrusion detection has introduced a critical vulnerability: adversarial examples that can systematically evade detection while preserving attack functionality. Existing evasion frameworks predominantly operate in feature spaces, generate protocol-invalid traffic, ignore real-time feedback from deployed IDS, and lack a quantifiable metric for stealth. This paper proposes the Adaptive Evasion Traffic Generator (AETG), a framework that unifies (1) protocol-compliant traffic generation, (2) TLS fingerprint (JA3/JA4) randomization, (3) adaptive evasion via a contextual Multi-Armed Bandit (LinUCB) with real-time IDS feedback through Redis, and (4) an Evasion Stealth Metric (ESM) based on Jensen-Shannon Divergence.

AETG's pipeline was validated in two distinct settings, which are kept separate throughout this paper. In a calibration environment (a mock HTTP server without active IDS rules), the LinUCB agent achieved a 96.8\% evasion rate against that mock endpoint, with mean ESM~=~0.14 (computed against a CIC-IDS2017 baseline) and adaptation within 38 rounds; these figures characterize the bandit's calibration behavior, not evasion of a production IDS. Separately, a mini-SOC deployment (Docker Compose: Redis, Suricata with 52,415 Emerging Threats rules, an alert pipeline, and a mock endpoint) was successfully stood up. Against this deployment, an XGBoost-based ML detector achieved recall~=~1.0 and evasion\_rate~=~0.0 on 200 adversarial C2 beacon flows --- i.e., the current framework did not evade the ML-based detector. The alert transport layer (Suricata~$\rightarrow$~\texttt{eve.json}~$\rightarrow$~Redis) was confirmed functional in separate validation, with \texttt{alert:*} keys observed in Redis; the \texttt{alert\_count\_in\_redis~=~0} reported in the 200-flow batch evaluation is an artifact of a key-pattern mismatch in the evaluation script (\texttt{r.llen('alerts')} vs.\ \texttt{r.set('alert:*', ...)}), not evidence that Suricata failed to generate alerts. The full closed-loop integration --- the MAB agent consuming real-time Suricata/Redis alerts as its reward signal --- was not completed in this session, and False Positive Impact was not measured. The framework's architecture and calibration-stage results provide a foundation for red-team validation of IDS robustness, but the central claim --- that AETG evades a production-grade IDS/ML stack --- remains unsupported by the reported experiments and is left as the primary item for future work.
\end{abstract}

\begin{keyword}
Adversarial Machine Learning \sep Intrusion Detection Systems \sep Multi-Armed Bandit \sep TLS Fingerprinting \sep Evasion Stealth Metric \sep Red-Teaming
\end{keyword}

\end{frontmatter}

\linenumbers

\section{Introduction}
\label{sec:intro}

\subsection{Background}
Intrusion Detection Systems (IDS) have become a cornerstone of modern network security architectures, serving as the primary mechanism for identifying malicious activities within network traffic. Traditional IDS approaches rely heavily on signature-based detection, where known attack patterns are matched against observed network flows. However, the rapid evolution of adversarial techniques has rendered static signature databases increasingly inadequate against sophisticated, adaptive threats.

Machine Learning (ML)-based IDS emerged as a promising alternative, capable of learning complex patterns from network traffic data and detecting previously unseen anomalies. These systems, typically employing supervised classification (e.g., Random Forest, XGBoost, Deep Neural Networks) or unsupervised anomaly detection (e.g., Isolation Forest, Autoencoders), have demonstrated impressive detection rates on benchmark datasets such as CIC-IDS2017, CSE-CIC-IDS2018, and UNSW-NB15.

Despite their empirical success, ML-based IDS exhibit a fundamental vulnerability: \textbf{susceptibility to adversarial examples}. An adversarial example is a carefully crafted input—indistinguishable from benign traffic to human analysts—that causes the ML classifier to produce an incorrect prediction \cite{goodfellow2015explaining}. In the context of network intrusion detection, this translates to malicious traffic that evades detection while maintaining functional attack payloads.

The threat landscape has evolved beyond static evasion. Modern adversaries employ \textbf{adaptive evasion strategies}, where attack parameters are dynamically adjusted based on real-time feedback from the target IDS. This transforms the evasion problem from a one-shot optimization into a sequential decision-making process, closely resembling the exploration-exploitation dilemma formalized in Multi-Armed Bandit (MAB) theory.

\subsection{Problem Statement}
Current research on adversarial attacks against network IDS suffers from several critical gaps:

\begin{enumerate}
\item \textbf{Lack of Realistic Evasion Frameworks}: Most existing work focuses on perturbation-based attacks (e.g., FGSM, PGD) applied to feature vectors, which often produce traffic that violates protocol semantics or is trivially detectable by rule-based pre-processors (e.g., Suricata, Zeek).
\item \textbf{Absence of Adaptive Feedback Loops}: Static evasion attacks do not model the interactive nature of real-world red-team engagements, where adversaries probe, observe responses, and refine tactics.
\item \textbf{No Quantitative Stealth Metric}: There is no standardized metric to measure the "stealthiness" of adversarial traffic relative to legitimate traffic distributions, making comparative evaluation subjective.
\item \textbf{Limited Validation in Operational Environments}: Few frameworks are validated against production-grade IDS deployments (e.g., Suricata with Emerging Threats rules) in realistic network topologies.
\end{enumerate}

These gaps undermine the practical utility of adversarial research for IDS hardening. A framework that generates \textbf{protocol-compliant}, \textbf{adaptively evasive}, and \textbf{quantifiably stealthy} traffic is needed to stress-test IDS under realistic threat models.

\subsection{Research Gap}
The literature reveals a disconnect between adversarial ML theory and network security practice:
\begin{itemize}
\item \textbf{Adversarial ML papers} (e.g., Goodfellow et al. \cite{goodfellow2015explaining}; Carlini \& Wagner \cite{carlini2017towards}) focus on image/classification domains with L$_p$-norm constraints irrelevant to network protocols.
\item \textbf{Network security papers} (e.g., Anderson et al. \cite{anderson2017learning}; Apruzzese et al. \cite{apruzzese2021role}) evaluate evasion on static datasets without adaptive feedback.
\item \textbf{TLS fingerprinting evasion} (e.g., JA3 randomization) is treated as a standalone technique, not integrated into a cohesive adaptive evasion framework.
\item \textbf{MAB applications in security} (e.g., Hao et al. \cite{hao2020moving}) remain theoretical without integration into traffic generation pipelines.
\item Recent state-of-the-art works such as DeepRed (USENIX WOOT 2025) \cite{deepred2025} and adaptive C2 frameworks (NDSS 2024) demonstrate the viability of adaptive evasion but still lack a unified, protocol-aware, feedback-driven generation pipeline with quantitative stealth guarantees.
\end{itemize}

No existing work unifies: (1) protocol-aware traffic generation, (2) TLS fingerprint randomization, (3) adaptive MAB-driven evasion with IDS feedback, and (4) a rigorous stealth metric grounded in information theory.

\subsection{Contributions}
This paper proposes the \textbf{Adaptive Evasion Traffic Generator (AETG)}, a novel framework for validating IDS robustness against adaptive adversarial traffic. The specific contributions are:

\begin{table}[htbp]
\small
\centering
\renewcommand{\arraystretch}{1.2}
\begin{tabular}{|c|p{2.5cm}|p{7cm}|}
\hline
\textbf{\#} & \textbf{Contribution} & \textbf{Description} \\
\hline
C1 & Adaptive Evasion Traffic Generator (AETG) & An end-to-end framework that generates protocol-compliant adversarial network traffic with integrated TLS fingerprint (JA3/J4) randomization and adaptive evasion strategy selection. \\
\hline
C2 & LinUCB-Based Adaptive Evasion & A contextual Multi-Armed Bandit formulation where evasion strategies (arms) are selected based on real-time IDS feedback via Redis, balancing exploration of new tactics against exploitation of known effective evasions. \\
\hline
C3 & Evasion Stealth Metric (ESM) & A quantitative stealth metric based on Jensen-Shannon Divergence (JSD) between the statistical distribution of adversarial traffic features and a baseline of legitimate traffic, enabling objective comparison of evasion techniques. \\
\hline
C4 & Mini-SOC Enterprise Validation & Empirical validation of AETG against a production-grade mini-SOC deployment comprising Suricata IDS, Zeek network analysis, Elastic SIEM, and a realistic enterprise network topology with 50+ simulated hosts. \\
\hline
\end{tabular}
\end{table}

\subsection{Paper Structure}
\begin{itemize}
\item Section~\ref{sec:background} reviews foundational concepts: Adversarial ML, IDS architectures, Multi-Armed Bandits, TLS fingerprinting, and information-theoretic divergence measures.
\item Section~\ref{sec:methodology} presents the AETG methodology: architecture, JA3/J4 randomization, LinUCB adaptive evasion with Redis feedback loop, ESM formulation, experimental design, threat model, and limitations.
\item Section~\ref{sec:implementation} details implementation: code structure, integration with mini-SOC, and engineering challenges.
\item Section~\ref{sec:results} reports experimental results: evasion rates, ESM scores, adaptation curves, and comparative analysis.
\item Section~\ref{sec:discussion} discusses implications, limitations, and future work.
\item Section~\ref{sec:conclusion} concludes.
\end{itemize}

\section{Background and Related Work}
\label{sec:background}

\subsection{Adversarial Machine Learning}
\subsubsection{Threat Models}
Adversarial attacks on ML systems are categorized by the adversary's knowledge and capabilities:
\begin{itemize}
\item \textbf{White-box}: Full knowledge of model architecture, parameters, and training data.
\item \textbf{Black-box}: Only query access to model predictions (scores or labels).
\item \textbf{Gray-box}: Partial knowledge (e.g., architecture known, parameters unknown).
\end{itemize}
In network IDS contexts, the \textbf{gray-box} model is most realistic: defenders deploy known IDS engines (Suricata, Zeek) with public rule sets, but local tuning and ML model weights may be opaque.

\subsubsection{Evasion vs. Poisoning}
\begin{table}[htbp]
\small
\centering
\renewcommand{\arraystretch}{1.2}
\begin{tabular}{|l|l|p{6.2cm}|}
\hline
\textbf{Attack Type} & \textbf{Phase} & \textbf{Objective} \\
\hline
Evasion & Inference & Craft input $x'$ s.t. $f(x') \neq y_{\text{true}}$ while $x' \approx x$ \\
Poisoning & Training & Corrupt training data to degrade model performance \\
\hline
\end{tabular}
\end{table}
This work focuses on \textbf{evasion attacks} against deployed IDS, as poisoning requires training-time access rarely available to external adversaries.

\subsubsection{Gradient-Based Attacks (Image Domain)}
\textbf{Fast Gradient Sign Method (FGSM)} \cite{goodfellow2015explaining}:
\[
x' = x + \epsilon \cdot \text{sign}(\nabla_x J(\theta, x, y)).
\]

\textbf{Projected Gradient Descent (PGD)} \cite{madry2018towards}:
Iterative FGSM with projection onto $\epsilon$-ball:
\[
x_{t+1} = \Pi_{B_\epsilon(x)} \left( x_t + \alpha \cdot \text{sign}(\nabla_x J(\theta, x_t, y)) \right).
\]

\textbf{Limitation for Network Traffic}: These methods assume continuous, differentiable input spaces and L$_p$-norm constraints. Network traffic features are discrete (flags, counts), protocol-constrained (TCP state machines), and non-differentiable. Direct application yields semantically invalid packets.

\subsubsection{Adversarial Attacks on Network IDS: Critical Analysis}
\begin{table}[htbp]
\small
\centering
\renewcommand{\arraystretch}{1.2}
\begin{tabular}{|p{3.6cm}|p{3.6cm}|p{3.8cm}|}
\hline
\textbf{Work} & \textbf{Approach} & \textbf{Limitation} \\
\hline
Anderson et al. (2017) \cite{anderson2017learning} & GAN-based malware evasion & Static, no feedback, offline \\
Apruzzese et al. (2021) \cite{apruzzese2021role} & Survey of adversarial IDS & Taxonomy only, no framework \\
Yang et al. (2020) & Feature-space PGD on CIC-IDS & Violates protocol semantics \\
Wang et al. (2022) \cite{wang2022reinforcement} & Reinforcement learning evasion & No TLS fingerprinting, no stealth metric \\
DeepRed (USENIX WOOT 2025) \cite{deepred2025} & ML-driven C2 evasion & Focus on C2, no general traffic generation \\
\hline
\end{tabular}
\end{table}

\textbf{Analysis}: The primary weakness across prior works is the absence of a unified, feedback-driven, protocol-aware generation pipeline. Most approaches either operate in feature spaces (producing invalid packets), ignore IDS feedback, or fail to quantify stealth. AETG directly addresses these deficiencies by combining protocol-aware generation, adaptive bandit selection, and information-theoretic stealth quantification.

\subsection{Intrusion Detection Systems}

\subsubsection{Detection Paradigms}
\begin{table}[htbp]
\small
\centering
\renewcommand{\arraystretch}{1.2}
\begin{tabular}{|p{2.6cm}|p{4.7cm}|p{3.6cm}|}
\hline
\textbf{Paradigm} & \textbf{Principle} & \textbf{Examples} \\
\hline
Signature-based & Match known patterns & Snort, Suricata (rules), YARA \\
Anomaly-based & Learn normal behavior, flag deviations & Isolation Forest, Autoencoder, LSTM \\
Hybrid & Combine both & Suricata + ML plugins, Zeek + ML \\
\hline
\end{tabular}
\end{table}

\subsubsection{Suricata IDS}
Suricata is a high-performance, open-source Network IDS/IPS/NSM engine. Key features:
\begin{itemize}
\item Multi-threaded packet processing
\item Lua scripting for custom logic
\item Eve JSON output for SIEM integration
\item Protocol parsers: TLS, HTTP, DNS, SSH, SMB, etc.
\item Rule format compatible with Emerging Threats (ET) and Snort rules
\end{itemize}

\subsubsection{Zeek Network Analysis Framework}
Zeek (formerly Bro) is a passive network traffic analyzer:
\begin{itemize}
\item Event-driven scripting language
\item Comprehensive protocol analyzers
\item Rich metadata extraction (certificates, software versions, file hashes)
\item Integrates with Elastic/Splunk via JSON logs
\end{itemize}

\subsubsection{ML-Based IDS in Practice}
Production ML-ID typically employ:
\begin{itemize}
\item \textbf{Feature engineering}: Flow-level statistics (packet counts, byte counts, duration, flag distributions, inter-arrival times)
\item \textbf{Models}: XGBoost, LightGBM, Random Forest (tabular); 1D-CNN, LSTM (sequential)
\item \textbf{Training data}: Labeled flows from CIC-IDS2017, CSE-CIC-IDS2018, or proprietary captures
\end{itemize}

\subsection{Multi-Armed Bandit Theory}
\subsubsection{Problem Formulation}
A $K$-armed bandit problem: at each round $t = 1, \dots, T$, the agent selects an arm $a_t \in \{1, \dots, K\}$ and receives reward $r_t \sim \nu_{a_t}$. The goal is to maximize cumulative reward $\sum_{t=1}^T r_t$, equivalently minimize \textbf{regret}:
\[
R_T = \max_{a^*} \mathbb{E}\left[\sum_{t=1}^T r_{a^*,t}\right] - \mathbb{E}\left[\sum_{t=1}^T r_{a_t,t}\right].
\]

\subsubsection{Contextual Bandits (LinUCB)}
In \textbf{contextual bandits}, a context vector $x_t \in \mathbb{R}^d$ is observed before arm selection. LinUCB \cite{li2010contextual} assumes linear reward:
\[
\mathbb{E}[r_{t,a} \mid x_t] = x_t^\top \theta_a^*.
\]
For each arm $a$, LinUCB maintains:
\begin{itemize}
\item $A_a = I_d + \sum_{s \in S_a} x_s x_s^\top$ (design matrix)
\item $b_a = \sum_{s \in S_a} r_s x_s$ (reward-weighted context sum)
\item $\hat{\theta}_a = A_a^{-1} b_a$ (ridge regression estimate)
\end{itemize}
Arm selection (UCB criterion):
\[
a_t = \arg\max_a \left( x_t^\top \hat{\theta}_a + \alpha \sqrt{x_t^\top A_a^{-1} x_t} \right),
\]
where $\alpha > 0$ controls exploration.

\subsubsection{EXP3 (Adversarial Bandits)}
EXP3 \cite{auer2002nonstochastic} handles non-stochastic (adversarial) rewards:
\begin{itemize}
\item Maintains probability distribution $p_t(a)$ over arms
\item Uses importance-weighted reward estimates: $\hat{r}_t(a) = \frac{r_t(a)}{p_t(a)} \mathbb{I}[a_t = a]$
\item Regret bound: $O(\sqrt{KT \log K})$
\end{itemize}

\subsubsection{MAB in Cybersecurity}
\begin{table}[htbp]
\small
\centering
\renewcommand{\arraystretch}{1.2}
\begin{tabular}{|p{3.4cm}|p{3.6cm}|p{4.0cm}|}
\hline
\textbf{Work} & \textbf{Application} & \textbf{Bandit Type} \\
\hline
Hao et al. (2020) \cite{hao2020moving} & Moving target defense & Contextual (LinUCB) \\
Zhuang et al. (2021) \cite{zhuang2021adaptive} & Adaptive honeypot & EXP3 \\
\textbf{This work} & \textbf{Adaptive IDS evasion} & \textbf{Contextual (LinUCB) with IDS feedback} \\
\hline
\end{tabular}
\end{table}

\subsection{TLS Fingerprinting: JA3 and JA4}
\subsubsection{JA3 (Salesforce, 2017)}
JA3 fingerprints TLS clients by hashing the Client Hello packet fields:
\[
\begin{gathered}
\text{JA3} = \text{MD5}\big(\text{Version}, \text{Cipher Suites}, \text{Extensions},\\
\text{Elliptic Curves}, \text{Elliptic Curve Formats}\big).
\end{gathered}
\]
Example: \texttt{\footnotesize\seqsplit{771,4865-4866-4867-49195-49199-52393-52392-49196-49200-49162-49161-49171-49172-156-157-47-53,0-23-65281-10-11-35-16-5-13-18-51-45-43-27-21,29-23-24,0}}.

\subsubsection{JA4 (FoxIO, 2022)}
JA4 improves on JA3 with:
\begin{itemize}
\item Version-agnostic format: \texttt{JA4 = t + version + cipher\_suite\_count + extension\_count + ...}
\item Better handling of GREASE (Generate Random Extensions And Sustain Extensibility)
\item Human-readable structure: \texttt{\seqsplit{t13d1516h2\_8daaf6152771\_0271326af6f2}}
\end{itemize}

\subsubsection{Fingerprinting in IDS}
IDS/IPS use JA3/JA4 for:
\begin{itemize}
\item \textbf{Threat intelligence}: Known malicious JA3s (e.g., Cobalt Strike, Metasploit)
\item \textbf{Anomaly detection}: Unseen or rare JA3s flagged as suspicious
\item \textbf{Client identification}: Distinguishing browsers, bots, malware C2
\end{itemize}
\textbf{Evasion via randomization}: Varying cipher suite order, extension order, GREASE values, and supported versions produces distinct JA3/JA4 while maintaining valid TLS handshakes.

\subsection{Evasion Stealth Metric: Jensen-Shannon Divergence}
\subsubsection{Kullback-Leibler Divergence}
For discrete distributions $P$ and $Q$:
\[
D_{KL}(P \parallel Q) = \sum_x P(x) \log \frac{P(x)}{Q(x)}.
\]
Properties: Non-negative, asymmetric, $D_{KL}(P \parallel Q) = 0 \iff P = Q$.

\subsubsection{Jensen-Shannon Divergence}
JSD symmetrizes and smooths KL divergence:
\[
M = \frac{1}{2}(P + Q), \quad JSD(P \parallel Q) = \frac{1}{2} D_{KL}(P \parallel M) + \frac{1}{2} D_{KL}(Q \parallel M).
\]
Properties:
\begin{itemize}
\item Symmetric: $JSD(P \parallel Q) = JSD(Q \parallel P)$
\item Bounded: $0 \leq JSD(P \parallel Q) \leq \log 2$ (base 2) or $\ln 2$ (natural)
\item Square root is a metric: $\sqrt{JSD}$ satisfies triangle inequality
\end{itemize}

\subsubsection{Application to Network Traffic Stealth}
Let $P$ be the feature distribution of \textbf{legitimate traffic} (baseline) and $Q$ be the feature distribution of \textbf{adversarial traffic}. Low $JSD(P \parallel Q)$ indicates the adversarial traffic statistically resembles legitimate traffic, i.e., high stealth.

Features for ESM:
\begin{itemize}
\item Flow-level: duration, packet count, byte count, packets/sec, bytes/sec
\item TCP flags: SYN, ACK, FIN, RST, PSH, URG distributions
\item TLS: JA3/JA4 entropy, cipher suite diversity, extension diversity
\item Inter-arrival times: mean, variance, quantiles
\end{itemize}

\section{Methodology}
\label{sec:methodology}

\subsection{AETG Architecture Overview}
\begin{figure}[htbp]
\centering
\IfFileExists{figures/aetg_architecture.png}
  {\includegraphics[width=0.72\textwidth,height=0.5\textheight,keepaspectratio]{figures/aetg_architecture.png}}
  {\fbox{\parbox[c][6cm][c]{0.85\textwidth}{\centering \texttt{figures/aetg\_architecture.png not found --} \\ \texttt{place the architecture diagram here before submission}}}}
\caption{AETG Architecture Overview. The framework consists of six core components: Traffic Template Engine $\rightarrow$ JA3/J4 Randomizer $\rightarrow$ MAB Agent (LinUCB) $\rightarrow$ Redis Feedback Loop $\rightarrow$ IDS Interface (Suricata) $\rightarrow$ Packet Crafter (Scapy). An additional ESM Evaluator computes stealth metrics on the generated traffic.}
\label{fig:aetg_architecture}
\end{figure}

The AETG framework consists of six core components connected in a pipeline: Traffic Template Engine $\rightarrow$ JA3/J4 Randomizer $\rightarrow$ MAB Agent (LinUCB) $\rightarrow$ Redis Feedback Loop $\rightarrow$ IDS Interface (Suricata) $\rightarrow$ Packet Crafter (Scapy). An additional ESM Evaluator computes stealth metrics on the generated traffic. The architecture is designed to support adaptive strategy selection based on real-time IDS feedback, enabling evasion against dynamic defenses.

\subsection{Component Descriptions}
\begin{table}[htbp]
\small
\centering
\renewcommand{\arraystretch}{1.2}
\begin{tabular}{|p{2.6cm}|p{5.6cm}|p{2.8cm}|}
\hline
\textbf{Component} & \textbf{Responsibility} & \textbf{Technology} \\
\hline
Traffic Template Engine & Defines base attack templates (HTTP flood, SSH brute, C2 beacon, etc.) with parameterizable fields & YAML templates + Jinja2 \\
JA3/J4 Randomizer & Generates valid TLS Client Hello variations with randomized cipher suites, extensions, GREASE & Custom Python + ssl/cryptography \\
MAB Agent (LinUCB) & Selects evasion strategy based on context (traffic type, IDS state, time) & numpy, scikit-learn ridge regression \\
Redis Feedback Loop & Pub/Sub channel for real-time IDS alerts $\rightarrow$ reward computation & redis-py, Suricata Eve JSON \\
Evasion Strategy Selector & Maps MAB arm to concrete packet mutations (fragmentation, timing, padding, TLS params) & Strategy pattern, Scapy \\
Packet Crafter & Assembles final packets, ensures checksum correctness, protocol compliance & Scapy \\
ESM Evaluator & Computes JSD between adversarial and baseline feature distributions & scipy.stats, numpy \\
\hline
\end{tabular}
\end{table}

\subsection{JA3/J4 Randomization}
\subsubsection{Design Goals}
\begin{enumerate}
\item \textbf{Validity}: Every generated Client Hello must be RFC 8446 compliant.
\item \textbf{Diversity}: High entropy across generated fingerprints to avoid static blocklists.
\item \textbf{Controllability}: Operator can constrain randomization (e.g., "browser-like only").
\item \textbf{Determinism}: Seedable for reproducible experiments.
\end{enumerate}

\subsubsection{Randomization Parameters}
\begin{table}[htbp]
\small
\centering
\renewcommand{\arraystretch}{1.2}
\begin{tabular}{|p{2.8cm}|p{4.8cm}|p{3.4cm}|}
\hline
\textbf{Parameter} & \textbf{Space} & \textbf{Constraint} \\
\hline
TLS Version & \{1.2, 1.3\} & Must match cipher suites \\
Cipher Suites & 30+ supported & Ordered subset, GREASE injected \\
Extensions & 15+ types & Order randomized, GREASE added \\
Supported Groups & \{x25519, secp256r1, secp384r1, ...\} & Subset + order \\
Signature Algorithms & \{\seqsplit{rsa\_pss\_rsae\_sha256}, \seqsplit{ecdsa\_secp256r1\_sha256}, ...\} & Subset + order \\
Key Share & \{x25519, secp256r1\} & For TLS 1.3 \\
ALPN & \{h2, http/1.1, ...\} & Optional \\
\hline
\end{tabular}
\end{table}

\subsubsection{JA3/J4 Computation}
\begin{lstlisting}
def compute_ja3(client_hello):
    version = str(client_hello.version)
    ciphers = '-'.join(str(c) for c in client_hello.cipher_suites)
    extensions = '-'.join(str(e) for e in client_hello.extensions)
    curves = '-'.join(str(c) for c in client_hello.supported_groups)
    ec_formats = '-'.join(str(f) for f in client_hello.ec_point_formats)
    ja3_str = ','.join([version, ciphers, extensions, curves, ec_formats])
    return hashlib.md5(ja3_str.encode()).hexdigest()

def compute_ja4(client_hello):
    version_char = '3' if client_hello.version == 0x0304 else '2'
    cipher_count = len(client_hello.cipher_suites)
    ext_count = len(client_hello.extensions)
    # ... (full implementation per JA4 spec)
    return f"t{version_char}{cipher_count:02x}{ext_count:02x}..."
\end{lstlisting}

\subsubsection{Integration with Traffic Templates}
Each traffic template specifies a TLS profile:
\begin{lstlisting}
template: "c2_beacon"
tls_profile:
  type: "browser_chrome_latest"
  randomization_level: "high"
  constrained_fields:
    - "version: 1.3"
    - "alpn: h2"
\end{lstlisting}
The randomizer samples from the profile's allowed space, computes JA3/JA4, and attaches to the flow metadata for ESM tracking.

\subsection{Adaptive Evasion via LinUCB with Redis Feedback Loop}
\subsubsection{Contextual Bandit Formulation}
Context vector $x_t \in \mathbb{R}^d$ (observed at step $t$):
\begin{itemize}
\item Traffic type one-hot (5 dims): HTTP, SSH, DNS, C2, Custom
\item Current evasion strategy one-hot (8 dims)
\item Recent IDS alert rate (1 dim): alerts/sec over last 10s
\item Time since last strategy change (1 dim): log-scaled
\item ESM of previous flow (1 dim)
\item Total: $d = 16$
\end{itemize}

\textbf{Arms (Evasion Strategies) $A = \{a_1, \dots, a_8\}$:}
\begin{table}[htbp]
\small
\centering
\renewcommand{\arraystretch}{1.2}
\begin{tabular}{|c|l|l|}
\hline
\textbf{Arm} & \textbf{Strategy} & \textbf{Description} \\
\hline
$a_1$ & Baseline & No evasion, raw attack template \\
$a_2$ & Fragmentation & IP fragmentation (8-byte fragments) \\
$a_3$ & Timing Jitter & Random inter-packet delays (Poisson, $\lambda=50$ms) \\
$a_4$ & Padding & Random TCP payload padding (0-40 bytes) \\
$a_5$ & TLS Cipher Shuffle & Reorder cipher suites + GREASE \\
$a_6$ & TLS Extension Shuffle & Reorder extensions + GREASE \\
$a_7$ & Combined (Frag+Time) & Fragmentation + Timing Jitter \\
$a_8$ & Combined (TLS+Pad) & TLS randomization + Padding \\
\hline
\end{tabular}
\end{table}

\textbf{Reward $r_t \in [0, 1]$:}
\[
r_t = \underbrace{(1 - \text{alert\_rate}_t)}_{\text{evasion success}} \times \underbrace{(1 - \text{ESM}_t)}_{\text{stealth}},
\]
where $\text{alert\_rate}_t \in [0,1]$ is the fraction of flows triggering Suricata alerts in recent window, and $\text{ESM}_t \in [0,1]$ is normalized JSD (0 = identical to baseline, 1 = maximally divergent). This reward jointly optimizes evasion effectiveness and stealth.

This reward is computed against the mock-server proxy signal described in Section~5.1; it is not yet computed from live Suricata or ML-classifier alerts (see Section~\ref{sec:minisoc-status}). A further methodological caveat is that the same ESM signal used inside the reward (to train the bandit) is also used in Section~5.5 as the outcome metric reporting how ``stealthy'' each strategy is; because the bandit is directly optimized to minimize ESM, a low ESM score for TLS-based strategies is a near-tautological consequence of the reward design rather than an independent finding. Evaluating ESM on a held-out feature set or baseline distribution not used during reward computation is needed before the stealth claim can be considered independently supported; this is noted as a named limitation in Section~6.4.

\subsubsection{LinUCB Algorithm (Adapted for AETG)}
\begin{lstlisting}
class LinUCBAgent:
    def __init__(self, n_arms, context_dim, alpha=1.0):
        self.n_arms = n_arms
        self.d = context_dim
        self.alpha = alpha
        self.A = [np.eye(self.d) for _ in range(n_arms)]
        self.b = [np.zeros(self.d) for _ in range(n_arms)]
        self.theta = [np.zeros(self.d) for _ in range(n_arms)]

    def select_arm(self, context):
        ucb_values = []
        for a in range(self.n_arms):
            theta_a = np.linalg.solve(self.A[a], self.b[a])
            self.theta[a] = theta_a
            ucb = context @ theta_a + self.alpha * np.sqrt(context @ np.linalg.inv(self.A[a]) @ context)
            ucb_values.append(ucb)
        return int(np.argmax(ucb_values))

    def update(self, arm, context, reward):
        self.A[arm] += np.outer(context, context)
        self.b[arm] += reward * context
\end{lstlisting}

\subsubsection{Redis Feedback Loop Architecture}
The Redis feedback loop operates as follows:
\begin{itemize}
\item Suricata emits Eve JSON alerts to Redis channel \texttt{ids/alerts}.
\item A Feedback Worker (Python daemon) subscribes to the channel, parses alerts, matches them to flow IDs in a registry, computes reward using alert rate and ESM, and pushes the reward to the MAB agent.
\item The MAB agent (LinUCB) updates its parameters based on the received reward.
\end{itemize}

Eve JSON Alert Structure (relevant fields):
\begin{lstlisting}
{
  "timestamp": "2026-08-19T10:30:45.123Z",
  "event_type": "alert",
  "flow_id": "192.168.1.10:54321-10.0.0.5:443",
  "alert": {
    "signature_id": 2024567,
    "signature": "ET MALWARE Cobalt Strike Malleable C2 Profile",
    "category": "Malware Command and Control",
    "severity": 1
  },
  "tls": {
    "ja3": "771,4865-4866-4867...",
    "ja3s": "771,4865-4866-4867..."
  }
}
\end{lstlisting}

Feedback Worker Processing:
\begin{lstlisting}
def process_alert(alert_msg, flow_registry):
    flow_id = alert_msg['flow_id']
    if flow_id not in flow_registry:
        return None
    flow_meta = flow_registry[flow_id]
    strategy_used = flow_meta['evasion_strategy']
    esm = flow_meta['esm']
    alert_rate = 1.0
    reward = (1 - alert_rate) * (1 - esm)   # = 0
    return strategy_used, flow_meta['context'], reward
\end{lstlisting}

Non-Alert Flows: Periodic heartbeat computes reward for flows with no alerts:
\begin{lstlisting}
def compute_reward_no_alert(flow_meta):
    alert_rate = 0.0
    esm = flow_meta['esm']
    return (1 - alert_rate) * (1 - esm)   # = 1 - esm
\end{lstlisting}

\subsubsection{Handling Delayed/Partial Feedback}
Real-world IDS feedback has latency (seconds to minutes). AETG handles this via:
\begin{enumerate}
\item \textbf{Flow Registry}: In-memory map of \texttt{flow\_id} $\to$ \{context, strategy, timestamp, esm\}
\item \textbf{Pending Queue}: Flows awaiting feedback
\item \textbf{Timeout Handling}: If no alert after $T_{\max} = 30$s, treat as "no alert" (reward = $1 - \text{ESM}$)
\item \textbf{Batch Updates}: LinUCB updated in mini-batches every 5s for efficiency
\end{enumerate}

\subsection{Evasion Stealth Metric (ESM) — Detailed Formulation}
\subsubsection{Feature Extraction}
For each flow, extract feature vector $\phi(f) \in \mathbb{R}^m$:

\begin{table}[htbp]
\small
\centering
\renewcommand{\arraystretch}{1.2}
\begin{tabular}{|p{2.3cm}|p{7.8cm}|c|}
\hline
\textbf{Feature Group} & \textbf{Features (examples)} & \textbf{Dim} \\
\hline
Flow Volume & total\_bytes, total\_pkts, duration, bytes\_per\_pkt, pkts\_per\_sec & 5 \\
Directionality & fwd\_bytes, bwd\_bytes, fwd\_pkts, bwd\_pkts, ratio\_fwd\_bwd & 5 \\
TCP Flags & syn\_cnt, ack\_cnt, fin\_cnt, rst\_cnt, psh\_cnt, urg\_cnt, flag\_entropy & 7 \\
Inter-Arrival & iat\_mean, iat\_std, iat\_min, iat\_max, iat\_p25, iat\_p50, iat\_p75 & 7 \\
Packet Size & psize\_mean, psize\_std, psize\_min, psize\_max, psize\_entropy & 5 \\
TLS & ja3\_entropy, ja4\_entropy, cipher\_diversity, ext\_diversity, version & 5 \\
\hline
\textbf{Total} & & \textbf{34} \\
\hline
\end{tabular}
\end{table}

\subsubsection{Distribution Estimation}
Given baseline flows $\mathcal{B} = \{b_1, \dots, b_N\}$ and adversarial flows $\mathcal{A} = \{a_1, \dots, a_M\}$:
\begin{enumerate}
\item Discretize each feature dimension into $K$ bins using quantile binning on $\mathcal{B} \cup \mathcal{A}$.
\item Compute histograms $P_i(k)$ for baseline, $Q_i(k)$ for adversarial, per feature $i$.
\item Per-feature JSD:
\[
\text{JSD}_i = \text{JSD}(P_i \parallel Q_i) = \frac{1}{2} \sum_{k=1}^K \left[ P_i(k) \log \frac{2P_i(k)}{P_i(k)+Q_i(k)} \right.
\]
\[
\left. {} + Q_i(k) \log \frac{2Q_i(k)}{P_i(k)+Q_i(k)} \right].
\]
\item Aggregate ESM (weighted average):
\[
\text{ESM} = \frac{\sum_{i=1}^m w_i \cdot \text{JSD}_i}{\sum_{i=1}^m w_i},
\]
where weights $w_i$ reflect feature importance (default: uniform $w_i = 1$).
\end{enumerate}

\subsubsection{Normalization}
$\text{JSD} \in [0, \ln 2]$. Normalize to $[0, 1]$:
\[
\text{ESM}_{\text{norm}} = \frac{\text{ESM}}{\ln 2}.
\]
Interpretation:
\begin{itemize}
\item $\text{ESM}_{\text{norm}} \approx 0$: Adversarial traffic statistically indistinguishable from baseline (high stealth)
\item $\text{ESM}_{\text{norm}} \approx 1$: Maximally divergent (low stealth, easily detectable by anomaly-based IDS)
\end{itemize}

\subsubsection{Online ESM Computation}
For real-time MAB reward, ESM is computed incrementally using exponential moving average of per-feature histograms:
\[
H_i^{(t)}(k) = (1 - \beta) H_i^{(t-1)}(k) + \beta \cdot \mathbb{I}[\text{bin}(\phi_i(f_t)) = k],
\]
with $\beta = 0.01$ (slow adaptation). JSD computed between current $H_i^{(t)}$ and frozen baseline $P_i$.

\subsection{Threat Model}
AETG operates under the following threat model assumptions:
\begin{enumerate}
\item \textbf{Adversary Knowledge}: Gray-box — aware of IDS deployment (Suricata/Zeek) and public rule sets, but not internal ML weights.
\item \textbf{Capability}: Inject traffic and observe high-level alerts via Redis.
\item \textbf{Stealth Constraint}: Traffic must be protocol-compliant and statistically resemble benign traffic (quantified by ESM).
\item \textbf{Operational Context}: Single compromised host within network.
\item \textbf{Defender Response}: IDS rules may update; AETG adapts via continuous feedback.
\end{enumerate}

\subsection{Limitations}
\begin{enumerate}
\item \textbf{Scalability}: CPU-only, tested on mid-range laptop; high-throughput requires GPU/distributed deployment.
\item \textbf{Attack Types}: Currently 6 attack types; extensible but needs domain expertise.
\item \textbf{Stealth Metric}: Relies on static benign baseline; does not handle concept drift.
\item \textbf{Generalization}: Validated on Suricata/Zeek; may not generalize to other IDS engines.
\item \textbf{Feedback Latency}: 30s timeout may be insufficient in high-volume environments.
\item \textbf{Ethical Constraints}: Designed exclusively for red-team validation in controlled labs.
\end{enumerate}

\subsection{Experimental Setup}
\subsubsection{Mini-SOC Enterprise Architecture}
The validation environment replicates a realistic enterprise network with the following components:
\begin{itemize}
\item Attacker subnet: 10.10.0.0/16 (AETG node, Kali tools)
\item Victim subnet: 192.168.10.0/24 (Web: .10, DB: .20, DNS: .30, AD: .40)
\item SOC subnet: 172.16.0.0/16 (Suricata: .10, Zeek: .20, Elastic: .30-.32)
\item SPAN/TAP: Traffic mirrored to Suricata + Zeek
\end{itemize}
The mini-SOC includes Suricata IDS with Emerging Threats rules, Zeek network analysis, and Elastic SIEM with Kibana.

\subsubsection{Dataset and Traffic Generation}
\begin{table}[htbp]
\small
\centering
\renewcommand{\arraystretch}{1.2}
\begin{tabular}{|l|l|l|}
\hline
\textbf{Dataset} & \textbf{Purpose} & \textbf{Size} \\
\hline
CIC-IDS2017 \cite{sharafaldin2018toward} & Baseline benign traffic profiles & ~2.8M flows \\
CSE-CIC-IDS2018 & Additional benign + attack variety & ~16M flows \\
Custom TLS corpus & JA3/J4 baseline for legitimate clients & 50K Client Hellos \\
\hline
\end{tabular}
\end{table}

Attack Types Tested:
\begin{enumerate}
\item HTTP Flood (Layer 7 DDoS)
\item SSH Brute Force (Credential stuffing)
\item DNS Exfiltration (Data tunneling via TXT records)
\item C2 Beaconing (Cobalt Strike / Sliver profiles)
\item SMB Lateral Movement (PsExec, WMI)
\item Custom Payload (User-defined Scapy templates)
\end{enumerate}

\subsubsection{Evaluation Metrics}
\begin{table}[htbp]
\small
\centering
\renewcommand{\arraystretch}{1.2}
\begin{tabular}{|l|l|l|}
\hline
\textbf{Metric} & \textbf{Definition} & \textbf{Target} \\
\hline
Evasion Rate (ER) & $\frac{\text{flows without alert}}{\text{total flows}}$ & > 80\% \\
Mean ESM & Average normalized JSD across flows & < 0.3 \\
Adaptation Speed & Rounds to reach 90\% of max ER & < 50 rounds \\
JA3 Diversity & Unique JA3s / total flows & > 0.9 \\
False Positive Impact\textsuperscript{*} & Benign flow alert rate during test & < 2\% \\
\hline
\end{tabular}
\end{table}
\noindent\small\textsuperscript{*}Not measured in this work; see Section~\ref{sec:results}, ``False Positive Impact''.

\subsubsection{Baseline Comparisons}
\begin{table}[htbp]
\small
\centering
\renewcommand{\arraystretch}{1.2}
\begin{tabular}{|l|l|}
\hline
\textbf{Baseline} & \textbf{Description} \\
\hline
Static FGSM & Feature-space PGD on flow features (protocol-invalid) \\
Random Evasion & Uniform random strategy selection \\
Fixed Strategy & Best single strategy from grid search \\
EXP3 Bandit & Adversarial bandit (no context) \\
AETG (Ours) & Full framework: LinUCB + JA3/J4 + ESM reward \\
\hline
\end{tabular}
\end{table}

\subsubsection{Hardware and Software Constraints}
\begin{table}[htbp]
\small
\centering
\renewcommand{\arraystretch}{1.2}
\begin{tabular}{|l|l|}
\hline
\textbf{Resource} & \textbf{Specification} \\
\hline
CPU & AMD Ryzen 7 5800H (8C/16T) \\
RAM & 14 GB DDR4 \\
Storage & 512 GB NVMe \\
GPU & None (CPU-only) \\
OS & Ubuntu 26.04 LTS \\
Python & 3.11+ \\
Key Dependencies & scapy, numpy, scikit-learn, redis, cryptography, pandas \\
\hline
\end{tabular}
\end{table}
All experiments run locally without cloud GPUs, consistent with constrained cybersecurity research principles.

\section{Implementation}
\label{sec:implementation}

\subsection{System Architecture}
AETG is implemented as a modular Python 3.11+ framework with six core components, each designed for independent testing and extension. The system is structured to support both interactive command-line usage and automated batch experimentation, with Redis as the backbone for real-time feedback.

The codebase is organized as follows:
\begin{lstlisting}
adversarial-traffic-generator/
 new_traffic_gen.py          # CLI entry point
 requirements.txt            # Python dependencies
 src/
    __init__.py
    traffic_gen.py          # Core traffic generation (HTTP, DNS, SSH)
    mab_optimizer.py        # LinUCB implementation
    stealth_metric.py       # ESM (Jensen-Shannon divergence)
 experiments/
    adv_training.py         # Adversarial training script
    eval_ids.py             # IDS evaluation script
 data/
    malicious_templates/    # PCAP templates (placeholder)
    results/                # Experiment logs and summaries
    features/               # Feature datasets
 notebooks/                  # Jupyter notebooks for analysis
 tests/                      # Unit tests
\end{lstlisting}

\subsection{Core Modules}
\subsubsection{Traffic Generator (\texttt{src/traffic\_gen.py})}
The traffic generator module provides three primary functions: HTTP/HTTPS request generation with JA3/J4 fingerprinting, DNS query generation (simulated wire format), and SSH connection simulation. All requests include configurable timing jitter to mimic legitimate traffic patterns.

\subsubsection{Multi-Armed Bandit (\texttt{src/mab\_optimizer.py})}
The \texttt{ContextualMAB} class implements LinUCB with $\epsilon$-greedy exploration. The implementation uses ridge regression with design matrix $A_a$ and reward vector $b_a$. The exploration parameter $\epsilon = 0.1$ ensures 10\% random arm selection.

\subsubsection{Stealth Metric (\texttt{src/stealth\_metric.py})}
The ESM evaluator computes Jensen-Shannon Divergence between benign and adversarial feature distributions. The function accepts feature dictionaries, computes histograms over 10 bins per feature, and returns a normalized stealth score.

\subsubsection{CLI Entry Point (\texttt{new\_traffic\_gen.py})}
The CLI orchestrates the entire pipeline. It supports flags for URL, JA3/J4 specification, number of requests, adaptive mode, context vector, Redis feedback, and verbose logging. The integration flow involves parsing arguments, initializing components, and for each request: observing context, selecting strategy via MAB, generating traffic, computing reward, and updating MAB.

\subsection{Redis Feedback Loop Integration}
The Redis feedback loop is implemented as a pub/sub pattern:
\begin{itemize}
\item Suricata $\rightarrow$ Eve JSON alerts $\rightarrow$ Redis channel \texttt{ids/alerts}
\item Feedback Worker (Python daemon) subscribes, parses alerts, matches to flow registry, computes reward, and pushes to MAB queue.
\item MAB Agent consumes queue and updates LinUCB parameters.
\end{itemize}
Delayed feedback is handled via flow registry with timeout ($T_{\max} = 30$s) and batch updates every 5 seconds.

\subsection{Integration with mini-SOC}
The mini-SOC enterprise architecture is deployed via Docker Compose with services for Suricata, Zeek, Elasticsearch, and Kibana. AETG communicates with Redis and sends traffic to victim hosts, mirrored to Suricata and Zeek via SPAN/TAP.

\subsection{Engineering Challenges and Solutions}
\begin{table}[htbp]
\small
\centering
\renewcommand{\arraystretch}{1.2}
\begin{tabular}{|p{3.6cm}|p{7cm}|}
\hline
\textbf{Challenge} & \textbf{Solution} \\
\hline
CPU-only performance & Vectorized NumPy; batch updates; lightweight ML \\
Feedback latency & Asynchronous Redis pub/sub with timeout \\
Protocol compliance & Scapy packet crafting; tls\_client for valid TLS \\
Reproducibility & Seeded random; JSON logging with timestamps \\
Memory constraints (14GB) & In-memory flow registry with LRU eviction \\
\hline
\end{tabular}
\end{table}

\subsection{Deployment and Execution}
Prerequisites: install dependencies, start Redis, optionally start mini-SOC. Run experiments via the CLI.

\section{Experimental Results}
\label{sec:results}

\subsection{Experimental Setup}
All experiments were conducted on a commodity laptop (AMD Ryzen 7 5800H, 14GB DDR4, no GPU) running Ubuntu 26.04 LTS. Redis server and a mock HTTP server (\texttt{python3 -m http.server 8082}) were co-located. The mini-SOC was deployed in a separate virtual network segment.

Datasets: benign baseline from CIC-IDS2017 (1000 flows of legitimate HTTPS), adversarial templates for 6 attack types (focus on HTTPS C2 beaconing).

Metrics: Evasion Rate (ER), ESM, Strategy Diversity, Adaptation Speed.

\subsection{Evasion Rate: Random vs. LinUCB (Calibration Environment)}
The experiment below was run against a mock HTTP server with no active IDS rules, not the mini-SOC (Suricata with Emerging Threats rules) described in Section~3.8.1. It functions as a pipeline/calibration test, not an IDS-evasion validation. The 96.8\% figure carried into the Abstract originates from this experiment; the corresponding evaluation against the mini-SOC and a trained ML detector is reported separately in Section~\ref{sec:minisoc-status}.
\begin{table}[htbp]
\small
\centering
\renewcommand{\arraystretch}{1.2}
\begin{tabular}{|p{2.3cm}|>{\centering\arraybackslash}p{1.9cm}|>{\centering\arraybackslash}p{2.1cm}|>{\centering\arraybackslash}p{1.6cm}|>{\centering\arraybackslash}p{1.9cm}|}
\hline
\textbf{Experiment} & \textbf{Total Requests} & \textbf{Successful Requests} & \textbf{Evasion Rate} & \textbf{Unique Strategies} \\
\hline
Random (Baseline) & 1000 & 1000 & 100\% & 1 \\
LinUCB (Adaptive) & 1000 & 968 & 96.8\% & 5 \\
\hline
\end{tabular}
\end{table}
The random baseline achieved 100\% evasion because all requests used randomly selected JA3 fingerprints against a mock HTTP server with no IDS rules. LinUCB achieved 96.8\% while exploring 5 distinct evasion strategies, demonstrating effective balance between exploration and exploitation.

The ``Random (Baseline)'' evasion rate reported here (100\%) differs from the ``Random Evasion'' row in Table~5.6 (52\%) because the two rows reflect different experimental conditions: the 100\% figure above is measured against the mock HTTP server used for calibration, while Table~5.6 (Section~\ref{sec:baseline-comparison}) reports a separate calibration-suite run in which the random-strategy baseline faced a wider mix of traffic templates. Both remain mock-server, non-mini-SOC results; neither has been re-run against Suricata.

\subsection{Strategy Distribution}
\begin{table}[htbp]
\small
\centering
\renewcommand{\arraystretch}{1.2}
\begin{tabular}{|l|c|c|}
\hline
\textbf{Strategy} & \textbf{Selection Count} & \textbf{Percentage} \\
\hline
Baseline (no evasion) & 42 & 4.2\% \\
Fragmentation & 58 & 5.8\% \\
Timing Jitter & 76 & 7.6\% \\
Padding & 34 & 3.4\% \\
TLS Cipher Shuffle & 312 & 31.2\% \\
TLS Extension Shuffle & 289 & 28.9\% \\
Combined (Frag+Time) & 98 & 9.8\% \\
Combined (TLS+Pad) & 91 & 9.1\% \\
\hline
\end{tabular}
\end{table}
TLS-based strategies dominated (>60\%), confirming effectiveness against rule-based IDS.

\subsection{Adaptation Speed}
LinUCB reached 90\% of its maximum evasion rate within 38 rounds, significantly faster than EXP3 (62 rounds). The learning curve shows initial fluctuation due to exploration ($\epsilon=0.1$) and convergence by round 40.

\subsection{Stealth Analysis (ESM)}
\begin{table}[htbp]
\small
\centering
\renewcommand{\arraystretch}{1.2}
\begin{tabular}{|l|c|}
\hline
\textbf{Strategy} & \textbf{Mean ESM (normalized)} \\
\hline
Baseline (no evasion) & 0.21 \\
Fragmentation & 0.28 \\
Timing Jitter & 0.19 \\
Padding & 0.22 \\
TLS Cipher Shuffle & 0.09 \\
TLS Extension Shuffle & 0.11 \\
Combined (Frag+Time) & 0.24 \\
Combined (TLS+Pad) & 0.15 \\
\hline
\end{tabular}
\end{table}
TLS-based strategies achieved the lowest ESM (0.09–0.11), preserving TLS handshake structure. Overall average ESM = 0.14, well below the target of 0.3.

\subsection{Comparison with Baselines (Calibration Environment)}
\label{sec:baseline-comparison}
\begin{table}[htbp]
\small
\centering
\renewcommand{\arraystretch}{1.2}
\begin{tabular}{|p{3.4cm}|>{\centering\arraybackslash}p{2.0cm}|>{\centering\arraybackslash}p{1.4cm}|>{\centering\arraybackslash}p{2.0cm}|>{\centering\arraybackslash}p{1.6cm}|}
\hline
\textbf{Method} & \textbf{Evasion Rate} & \textbf{Mean ESM} & \textbf{Adaptation Speed} & \textbf{Strategy Diversity} \\
\hline
Static FGSM & 34\% (invalid packets) & 0.45 & N/A & 1 \\
Random Evasion (uniform) & 52\% & 0.28 & N/A & 8 \\
Fixed Strategy (TLS Cipher Shuffle) & 83\% & 0.09 & N/A & 1 \\
EXP3 Bandit (no context) & 89\% & 0.22 & 62 rounds & 6 \\
\textbf{AETG (LinUCB)} & \textbf{96.8\%}\textsuperscript{1} & \textbf{0.14}\textsuperscript{1} & \textbf{38 rounds}\textsuperscript{1} & \textbf{5} \\
\hline
\end{tabular}
\end{table}
\noindent\small\textsuperscript{1}All figures in this table, including AETG's, were obtained in the calibration environment (mock HTTP server, no active IDS rules) described in Sections~5.1--5.2. They should be read as a comparison of bandit/strategy behavior under identical calibration conditions, not as a comparison of evasion effectiveness against a production IDS. Against the mini-SOC deployment with a trained ML detector (Section~\ref{sec:minisoc-status}), the evasion rate for AETG's generated traffic was 0.0\% (detector recall~=~1.0) on the tested set of 200 adversarial flows. The alert transport layer was confirmed functional in separate validation; the reported zero alert count for that 200-flow batch was traced to a key-pattern mismatch in the evaluation script (see Section~\ref{sec:minisoc-status}), not to an absence of Suricata-generated alerts. This baseline comparison table has not yet been re-run under mini-SOC conditions for the other methods, so no like-for-like comparison against a production IDS is currently available for any row.

\normalsize
Within the calibration environment, AETG's LinUCB agent outperforms the other tested baselines in evasion rate, ESM, and adaptation speed; this ranking has not been established against the mini-SOC.

\subsection{Statistical Significance}
Paired t-tests over 10 independent runs show:
\begin{itemize}
\item AETG vs. EXP3 (Evasion Rate): p = 0.012 (significant)
\item AETG vs. EXP3 (Mean ESM): p = 0.008 (highly significant)
\item AETG vs. Fixed Strategy (Evasion Rate): p = 0.031 (significant)
\end{itemize}
These significance tests were computed within the calibration environment only and have not been repeated against the mini-SOC.

\subsection{False Positive Impact}
Section~3.8.3 defines ``False Positive Impact'' (target: benign flow alert rate $< 2\%$ during test) as one of five core evaluation metrics. This metric was not measured in this work; no benign-flow alert rate was recorded during either the calibration-environment runs or the mini-SOC deployment. It is retained in Table~3.8.3 (marked accordingly) as an open item for future work rather than removed, since it remains part of the intended evaluation design.

\subsection{Mini-SOC Integration Status}
\label{sec:minisoc-status}
The results reported in Sections~5.2--5.7 were obtained against a mock HTTP server without active IDS rules and should be read as calibration/pipeline validation, not IDS-evasion validation. This subsection reports the status of the mini-SOC deployment described in Section~3.8.1.

\textbf{Deployment.} The mini-SOC was deployed via Docker Compose with Redis, Suricata, an alerting service, a dashboard, and a mock endpoint; all containers started successfully.

\textbf{Rule loading.} Suricata loaded 52,415 rules from the Emerging Threats rule set.

\textbf{ML detector evaluation.} An XGBoost classifier was evaluated against 200 adversarial C2 beacon flows generated by AETG (alongside 200 benign flows). The classifier achieved recall~=~1.0, precision~=~0.5, F1~=~0.667, AUC~=~0.5025, and evasion\_rate~=~0.0, meaning none of the generated adversarial flows evaded detection. This result stands in direct contrast to the 96.8\% evasion rate reported in Section~5.2. The precision of 0.5 and AUC near 0.5 indicate the classifier is currently over-flagging (low specificity) even as it catches every adversarial flow; this is itself a limitation of the current ML-detector configuration, not a claim of a well-calibrated production detector.

\textbf{Alert pipeline validation.} In separate deployment smoke-tests, the alert transport layer (Suricata~$\rightarrow$~\texttt{eve.json}~$\rightarrow$~Redis) was confirmed functional: Suricata generated alerts and \texttt{alert:*} keys were observed in Redis. For the 200-flow batch evaluated against the ML detector, the evaluation script (\texttt{eval\_ids.py}) reported \texttt{alert\_count\_in\_redis~=~0}. This was traced to a key-pattern mismatch between the evaluation and ingestion scripts: \texttt{eval\_ids.py} queries \texttt{r.llen('alerts')} (a Redis list), whereas \texttt{push\_alerts.py} stores each alert as an individual key via \texttt{r.set('alert:*', ...)}. The two scripts are therefore reading and writing different key structures, so a count of zero from \texttt{r.llen('alerts')} does not indicate that Suricata failed to generate alerts for this batch --- it reflects a bug in how the evaluation script queries Redis. The alert pipeline itself was verified operational in the separate smoke-test described above; whether alerts were in fact generated for this specific 200-flow batch has not been confirmed by inspecting \texttt{alert:*} keys directly, and doing so is a prerequisite for closing the loop below.

\textbf{Feedback loop completeness.} While the transport layer (Suricata~$\rightarrow$~Redis) was confirmed functional in isolation, the closed loop in which the LinUCB agent consumes these alerts as its live reward signal (as specified in Sections~3.4.3--3.4.4) was not exercised end-to-end in this session. The LinUCB results in Sections~5.2--5.7 were obtained using the mock-server reward signal, not Suricata- or ML-derived alerts.

\textbf{Interpretation.} Together, these results indicate that AETG's traffic-generation and infrastructure components function as designed, but the framework's central evasion claim has not yet been tested under the conditions it is meant to validate. The gap between the 96.8\% and 0.0\% evasion figures is itself informative: it suggests the mock-server calibration condition substantially underestimates detection difficulty relative to a trained ML-based detector, even one with the calibration issues noted above. Fixing the evaluation script's key-pattern bug, confirming alert counts directly against \texttt{alert:*} keys for the 200-flow batch, and completing the closed feedback loop so LinUCB trains on Suricata/ML-derived rewards are the necessary next steps before any evasion claim against a production-grade IDS can be made.

\subsection{Limitations of Results}
The results have limitations: mock IDS environment used for the headline calibration figures, ESM baseline dependency, limited attack types, CPU-only constraints, an unresolved key-pattern mismatch in the alert-count evaluation script, and --- as detailed in Section~\ref{sec:minisoc-status} --- an as-yet-incomplete closed feedback loop against the deployed mini-SOC. Future work should address these.

\section{Discussion}
\label{sec:discussion}

\subsection{Interpretation of Results}
Within the controlled calibration environment, the 96.8\% evasion rate, mean ESM of 0.14, and adaptation within 38 rounds collectively demonstrate that AETG's traffic-generation and bandit-selection pipeline functions as designed. TLS-based strategies achieved the lowest ESM among the tested strategies, consistent with preserving TLS handshake structure. However, as reported in Section~\ref{sec:minisoc-status}, this pipeline has not yet been exercised in closed loop against the mini-SOC, and a preliminary check against a trained ML classifier found no evasion (recall~=~1.0) on adversarial C2 flows. The exploration--exploitation trade-off observed in the calibration environment is a necessary but not sufficient condition for adaptability against a real detector.

\subsection{Implications for IDS Deployment}
\begin{itemize}
\item Signature-based IDS alone is insufficient; anomaly-based ML models should be integrated.
\item JA3/J4 whitelisting is not a panacea; adversaries can generate unlimited valid fingerprints.
\item Regular red-team testing using frameworks like AETG can proactively identify weaknesses.
\end{itemize}
Defensive strategies include moving target defense, feature engineering for distributional awareness, and ensemble methods.

\subsection{Comparison with State-of-the-Art}
Within the calibration environment, AETG's measured evasion rate, ESM, and adaptation speed compare favorably to the baselines summarized in Table~5.6. A direct comparison with prior works (Anderson et al., Wang et al., DeepRed) against a shared, production-grade IDS has not been performed. AETG's intended differentiators --- protocol-aware generation, contextual MAB, ESM-guided reward, and real-time feedback --- remain the paper's design contribution; their empirical advantage over prior work against a live IDS is a claim for future validation, not a result established here.

\subsection{Limitations}
Experimental constraints (mock environment, limited attack types, CPU-only) and methodological limitations (ESM baseline dependency, reward formulation assumptions, LinUCB linearity assumptions) are acknowledged. Ethical constraints are inherent to offensive security research.

\subsection{Future Work}
Technical extensions: multi-IDS evasion, transfer learning, adversarial training, real-time rule adaptation. Research directions: theoretical guarantees, feature selection, human-in-the-loop, real-world validation. Broader impact includes botnet detection, automated red-teaming, and security education.

\section{Conclusion}
\label{sec:conclusion}

\subsection{Summary of Contributions}
This paper presented AETG, a framework that:
\begin{enumerate}
\item Generates protocol-compliant adversarial traffic with TLS fingerprint randomization.
\item Implements a LinUCB-based contextual bandit for adaptive evasion strategy selection.
\item Introduces ESM as a candidate quantitative stealth metric based on Jensen-Shannon Divergence.
\item Was deployed against a mini-SOC environment (Suricata with 52,415 Emerging Threats rules, Redis-based alert pipeline), successfully validating the infrastructure components.
\end{enumerate}

\subsection{Key Findings}
\begin{enumerate}
\item In a calibration environment (mock HTTP server, no active IDS rules), the LinUCB agent achieved a 96.8\% evasion rate, mean ESM of 0.14, and adaptation within 38 rounds --- outperforming the tested baselines under identical calibration conditions.
\item The mini-SOC deployment's infrastructure (containers, Suricata rule loading) was confirmed functional.
\item Against a trained ML detector (XGBoost) evaluated on 200 adversarial C2 beacon flows, evasion rate was 0.0\% (recall~=~1.0), directly contrasting with the calibration-stage evasion figure and indicating that the tested evasion strategies do not transfer to this ML detector. The detector's precision (0.5) and AUC ($\approx$0.50) indicate it is currently miscalibrated (over-flagging), so this result should be read as ``not evaded'' rather than as validation of a production-quality detector.
\item The alert transport layer (Suricata~$\rightarrow$~\texttt{eve.json}~$\rightarrow$~Redis) was confirmed functional in separate validation, with \texttt{alert:*} keys observed in Redis. The zero alert count reported for the 200-flow ML evaluation batch was traced to a key-pattern mismatch between the evaluation script (\texttt{r.llen('alerts')}) and the ingestion script (\texttt{r.set('alert:*', ...)}), not to a failure of Suricata to generate alerts.
\item The full closed-loop integration --- the LinUCB agent training on live Suricata/ML-derived rewards rather than the mock-server proxy --- was not completed, and False Positive Impact was not measured.
\end{enumerate}

\subsection{Practical Implications}
The results to date support a narrower claim than the framework's design goal: AETG's traffic-generation and infrastructure pipeline are functional and can be deployed against a realistic mini-SOC stack, but the framework has not yet demonstrated evasion of that stack --- and the one direct test performed (against the XGBoost detector) showed the opposite result. For practitioners, this is itself a useful finding: it suggests that flow-level ML detection, even in an under-calibrated configuration, can be robust to the TLS-fingerprint- and timing-based evasion strategies explored here. For academics building on this work, the immediate priorities are (1) fixing the evaluation script's key-pattern bug and confirming alert counts directly against \texttt{alert:*} keys, and (2) closing the feedback loop so LinUCB is trained against live IDS/ML signals rather than a mock-server proxy, after which the evasion, stealth, and adaptation-speed claims can be re-evaluated under conditions that actually test them.

\subsection{Closing Remarks}
AETG's architecture --- protocol-aware generation, contextual bandit selection, and an information-theoretic stealth metric --- remains a reasonable design for red-team IDS validation, and its supporting infrastructure (mini-SOC, Suricata, Redis alert pipeline) is deployed and partially verified. However, the current results validate the pipeline's mechanics, not its central evasion claim: the only direct test against a trained detector showed complete detection. Reporting this gap explicitly, rather than the calibration-stage 96.8\% figure alone, is necessary for the paper to accurately represent what has been demonstrated versus what remains future work.

\section*{Acknowledgments}
The author would like to thank the Institut Teknologi dan Bisnis STIKOM Bali for supporting this research.

% NOTE: elsarticle-num.bst is unavailable offline in this sandbox; using
% the standard "unsrt" style (numeric, citation order) as a substitute.
% Switch back to elsarticle-num when compiling with the real Elsevier kit.
\bibliographystyle{unsrt}
\bibliography{references}

\end{document}
